%% bare_jrnl.tex
%% V1.4b
%% 2015/08/26
%% by Michael Shell
%% see http://www.michaelshell.org/
%% for current contact information.
%%
%% This is a skeleton file demonstrating the use of IEEEtran.cls
%% (requires IEEEtran.cls version 1.8b or later) with an IEEE
%% journal paper.
%%
%% Support sites:
%% http://www.michaelshell.org/tex/ieeetran/
%% http://www.ctan.org/pkg/ieeetran
%% and
%% http://www.ieee.org/

%%*************************************************************************
%% Legal Notice:
%% This code is offered as-is without any warranty either expressed or
%% implied; without even the implied warranty of MERCHANTABILITY or
%% FITNESS FOR A PARTICULAR PURPOSE! 
%% User assumes all risk.
%% In no event shall the IEEE or any contributor to this code be liable for
%% any damages or losses, including, but not limited to, incidental,
%% consequential, or any other damages, resulting from the use or misuse
%% of any information contained here.
%%
%% All comments are the opinions of their respective authors and are not
%% necessarily endorsed by the IEEE.
%%
%% This work is distributed under the LaTeX Project Public License (LPPL)
%% ( http://www.latex-project.org/ ) version 1.3, and may be freely used,
%% distributed and modified. A copy of the LPPL, version 1.3, is included
%% in the base LaTeX documentation of all distributions of LaTeX released
%% 2003/12/01 or later.
%% Retain all contribution notices and credits.
%% ** Modified files should be clearly indicated as such, including  **
%% ** renaming them and changing author support contact information. **
%%*************************************************************************


% *** Authors should verify (and, if needed, correct) their LaTeX system  ***
% *** with the testflow diagnostic prior to trusting their LaTeX platform ***
% *** with production work. The IEEE's font choices and paper sizes can   ***
% *** trigger bugs that do not appear when using other class files.       ***                          ***
% The testflow support page is at:
% http://www.michaelshell.org/tex/testflow/


\documentclass[journal]{IEEEtran}
%
% If IEEEtran.cls has not been installed into the LaTeX system files,
% manually specify the path to it like:
% \documentclass[journal]{../sty/IEEEtran}

\usepackage{algorithm}
\usepackage{algorithmic}
\usepackage{graphicx}
\usepackage{subcaption}
\usepackage{amsmath}
\usepackage{hyperref}
\usepackage{float}
\usepackage{xcolor}





% Some very useful LaTeX packages include:
% (uncomment the ones you want to load)


% *** MISC UTILITY PACKAGES ***
%
%\usepackage{ifpdf}
% Heiko Oberdiek's ifpdf.sty is very useful if you need conditional
% compilation based on whether the output is pdf or dvi.
% usage:
% \ifpdf
%   % pdf code
% \else
%   % dvi code
% \fi
% The latest version of ifpdf.sty can be obtained from:
% http://www.ctan.org/pkg/ifpdf
% Also, note that IEEEtran.cls V1.7 and later provides a builtin
% \ifCLASSINFOpdf conditional that works the same way.
% When switching from latex to pdflatex and vice-versa, the compiler may
% have to be run twice to clear warning/error messages.






% *** CITATION PACKAGES ***
%
%\usepackage{cite}
% cite.sty was written by Donald Arseneau
% V1.6 and later of IEEEtran pre-defines the format of the cite.sty package
% \cite{} output to follow that of the IEEE. Loading the cite package will
% result in citation numbers being automatically sorted and properly
% "compressed/ranged". e.g., [1], [9], [2], [7], [5], [6] without using
% cite.sty will become [1], [2], [5]--[7], [9] using cite.sty. cite.sty's
% \cite will automatically add leading space, if needed. Use cite.sty's
% noadjust option (cite.sty V3.8 and later) if you want to turn this off
% such as if a citation ever needs to be enclosed in parenthesis.
% cite.sty is already installed on most LaTeX systems. Be sure and use
% version 5.0 (2009-03-20) and later if using hyperref.sty.
% The latest version can be obtained at:
% http://www.ctan.org/pkg/cite
% The documentation is contained in the cite.sty file itself.






% *** GRAPHICS RELATED PACKAGES ***
%
\ifCLASSINFOpdf
  % \usepackage[pdftex]{graphicx}
  % declare the path(s) where your graphic files are
  % \graphicspath{{../pdf/}{../jpeg/}}
  % and their extensions so you won't have to specify these with
  % every instance of \includegraphics
  % \DeclareGraphicsExtensions{.pdf,.jpeg,.png}
\else
  % or other class option (dvipsone, dvipdf, if not using dvips). graphicx
  % will default to the driver specified in the system graphics.cfg if no
  % driver is specified.
  % \usepackage[dvips]{graphicx}
  % declare the path(s) where your graphic files are
  % \graphicspath{{../eps/}}
  % and their extensions so you won't have to specify these with
  % every instance of \includegraphics
  % \DeclareGraphicsExtensions{.eps}
\fi
% graphicx was written by David Carlisle and Sebastian Rahtz. It is
% required if you want graphics, photos, etc. graphicx.sty is already
% installed on most LaTeX systems. The latest version and documentation
% can be obtained at: 
% http://www.ctan.org/pkg/graphicx
% Another good source of documentation is "Using Imported Graphics in
% LaTeX2e" by Keith Reckdahl which can be found at:
% http://www.ctan.org/pkg/epslatex
%
% latex, and pdflatex in dvi mode, support graphics in encapsulated
% postscript (.eps) format. pdflatex in pdf mode supports graphics
% in .pdf, .jpeg, .png and .mps (metapost) formats. Users should ensure
% that all non-photo figures use a vector format (.eps, .pdf, .mps) and
% not a bitmapped formats (.jpeg, .png). The IEEE frowns on bitmapped formats
% which can result in "jaggedy"/blurry rendering of lines and letters as
% well as large increases in file sizes.
%
% You can find documentation about the pdfTeX application at:
% http://www.tug.org/applications/pdftex





% *** MATH PACKAGES ***
%
%\usepackage{amsmath}
% A popular package from the American Mathematical Society that provides
% many useful and powerful commands for dealing with mathematics.
%
% Note that the amsmath package sets \interdisplaylinepenalty to 10000
% thus preventing page breaks from occurring within multiline equations. Use:
%\interdisplaylinepenalty=2500
% after loading amsmath to restore such page breaks as IEEEtran.cls normally
% does. amsmath.sty is already installed on most LaTeX systems. The latest
% version and documentation can be obtained at:
% http://www.ctan.org/pkg/amsmath





% *** SPECIALIZED LIST PACKAGES ***
%
%\usepackage{algorithmic}
% algorithmic.sty was written by Peter Williams and Rogerio Brito.
% This package provides an algorithmic environment fo describing algorithms.
% You can use the algorithmic environment in-text or within a figure
% environment to provide for a floating algorithm. Do NOT use the algorithm
% floating environment provided by algorithm.sty (by the same authors) or
% algorithm2e.sty (by Christophe Fiorio) as the IEEE does not use dedicated
% algorithm float types and packages that provide these will not provide
% correct IEEE style captions. The latest version and documentation of
% algorithmic.sty can be obtained at:
% http://www.ctan.org/pkg/algorithms
% Also of interest may be the (relatively newer and more customizable)
% algorithmicx.sty package by Szasz Janos:
% http://www.ctan.org/pkg/algorithmicx




% *** ALIGNMENT PACKAGES ***
%
%\usepackage{array}
% Frank Mittelbach's and David Carlisle's array.sty patches and improves
% the standard LaTeX2e array and tabular environments to provide better
% appearance and additional user controls. As the default LaTeX2e table
% generation code is lacking to the point of almost being broken with
% respect to the quality of the end results, all users are strongly
% advised to use an enhanced (at the very least that provided by array.sty)
% set of table tools. array.sty is already installed on most systems. The
% latest version and documentation can be obtained at:
% http://www.ctan.org/pkg/array


% IEEEtran contains the IEEEeqnarray family of commands that can be used to
% generate multiline equations as well as matrices, tables, etc., of high
% quality.




% *** SUBFIGURE PACKAGES ***
%\ifCLASSOPTIONcompsoc
%  \usepackage[caption=false,font=normalsize,labelfont=sf,textfont=sf]{subfig}
%\else
%  \usepackage[caption=false,font=footnotesize]{subfig}
%\fi
% subfig.sty, written by Steven Douglas Cochran, is the modern replacement
% for subfigure.sty, the latter of which is no longer maintained and is
% incompatible with some LaTeX packages including fixltx2e. However,
% subfig.sty requires and automatically loads Axel Sommerfeldt's caption.sty
% which will override IEEEtran.cls' handling of captions and this will result
% in non-IEEE style figure/table captions. To prevent this problem, be sure
% and invoke subfig.sty's "caption=false" package option (available since
% subfig.sty version 1.3, 2005/06/28) as this is will preserve IEEEtran.cls
% handling of captions.
% Note that the Computer Society format requires a larger sans serif font
% than the serif footnote size font used in traditional IEEE formatting
% and thus the need to invoke different subfig.sty package options depending
% on whether compsoc mode has been enabled.
%
% The latest version and documentation of subfig.sty can be obtained at:
% http://www.ctan.org/pkg/subfig




% *** FLOAT PACKAGES ***
%
%\usepackage{fixltx2e}
% fixltx2e, the successor to the earlier fix2col.sty, was written by
% Frank Mittelbach and David Carlisle. This package corrects a few problems
% in the LaTeX2e kernel, the most notable of which is that in current
% LaTeX2e releases, the ordering of single and double column floats is not
% guaranteed to be preserved. Thus, an unpatched LaTeX2e can allow a
% single column figure to be placed prior to an earlier double column
% figure.
% Be aware that LaTeX2e kernels dated 2015 and later have fixltx2e.sty's
% corrections already built into the system in which case a warning will
% be issued if an attempt is made to load fixltx2e.sty as it is no longer
% needed.
% The latest version and documentation can be found at:
% http://www.ctan.org/pkg/fixltx2e


%\usepackage{stfloats}
% stfloats.sty was written by Sigitas Tolusis. This package gives LaTeX2e
% the ability to do double column floats at the bottom of the page as well
% as the top. (e.g., "\begin{figure*}[!b]" is not normally possible in
% LaTeX2e). It also provides a command:
%\fnbelowfloat
% to enable the placement of footnotes below bottom floats (the standard
% LaTeX2e kernel puts them above bottom floats). This is an invasive package
% which rewrites many portions of the LaTeX2e float routines. It may not work
% with other packages that modify the LaTeX2e float routines. The latest
% version and documentation can be obtained at:
% http://www.ctan.org/pkg/stfloats
% Do not use the stfloats baselinefloat ability as the IEEE does not allow
% \baselineskip to stretch. Authors submitting work to the IEEE should note
% that the IEEE rarely uses double column equations and that authors should try
% to avoid such use. Do not be tempted to use the cuted.sty or midfloat.sty
% packages (also by Sigitas Tolusis) as the IEEE does not format its papers in
% such ways.
% Do not attempt to use stfloats with fixltx2e as they are incompatible.
% Instead, use Morten Hogholm'a dblfloatfix which combines the features
% of both fixltx2e and stfloats:
%
% \usepackage{dblfloatfix}
% The latest version can be found at:
% http://www.ctan.org/pkg/dblfloatfix




%\ifCLASSOPTIONcaptionsoff
%  \usepackage[nomarkers]{endfloat}
% \let\MYoriglatexcaption\caption
% \renewcommand{\caption}[2][\relax]{\MYoriglatexcaption[#2]{#2}}
%\fi
% endfloat.sty was written by James Darrell McCauley, Jeff Goldberg and 
% Axel Sommerfeldt. This package may be useful when used in conjunction with 
% IEEEtran.cls'  captionsoff option. Some IEEE journals/societies require that
% submissions have lists of figures/tables at the end of the paper and that
% figures/tables without any captions are placed on a page by themselves at
% the end of the document. If needed, the draftcls IEEEtran class option or
% \CLASSINPUTbaselinestretch interface can be used to increase the line
% spacing as well. Be sure and use the nomarkers option of endfloat to
% prevent endfloat from "marking" where the figures would have been placed
% in the text. The two hack lines of code above are a slight modification of
% that suggested by in the endfloat docs (section 8.4.1) to ensure that
% the full captions always appear in the list of figures/tables - even if
% the user used the short optional argument of \caption[]{}.
% IEEE papers do not typically make use of \caption[]'s optional argument,
% so this should not be an issue. A similar trick can be used to disable
% captions of packages such as subfig.sty that lack options to turn off
% the subcaptions:
% For subfig.sty:
% \let\MYorigsubfloat\subfloat
% \renewcommand{\subfloat}[2][\relax]{\MYorigsubfloat[]{#2}}
% However, the above trick will not work if both optional arguments of
% the \subfloat command are used. Furthermore, there needs to be a
% description of each subfigure *somewhere* and endfloat does not add
% subfigure captions to its list of figures. Thus, the best approach is to
% avoid the use of subfigure captions (many IEEE journals avoid them anyway)
% and instead reference/explain all the subfigures within the main caption.
% The latest version of endfloat.sty and its documentation can obtained at:
% http://www.ctan.org/pkg/endfloat
%
% The IEEEtran \ifCLASSOPTIONcaptionsoff conditional can also be used
% later in the document, say, to conditionally put the References on a 
% page by themselves.




% *** PDF, URL AND HYPERLINK PACKAGES ***
%
%\usepackage{url}
% url.sty was written by Donald Arseneau. It provides better support for
% handling and breaking URLs. url.sty is already installed on most LaTeX
% systems. The latest version and documentation can be obtained at:
% http://www.ctan.org/pkg/url
% Basically, \url{my_url_here}.




% *** Do not adjust lengths that control margins, column widths, etc. ***
% *** Do not use packages that alter fonts (such as pslatex).         ***
% There should be no need to do such things with IEEEtran.cls V1.6 and later.
% (Unless specifically asked to do so by the journal or conference you plan
% to submit to, of course. )


% correct bad hyphenation here
\hyphenation{op-tical net-works semi-conduc-tor}


\begin{document}
%
% paper title
% Titles are generally capitalized except for words such as a, an, and, as,
% at, but, by, for, in, nor, of, on, or, the, to and up, which are usually
% not capitalized unless they are the first or last word of the title.
% Linebreaks \\ can be used within to get better formatting as desired.
% Do not put math or special symbols in the title.
\title{Introdução a Processamento de Imagens\newline Trabalho Final}
%
%
% author names and IEEE memberships
% note positions of commas and nonbreaking spaces ( ~ ) LaTeX will not break
% a structure at a ~ so this keeps an author's name from being broken across
% two lines.
% use \thanks{} to gain access to the first footnote area
% a separate \thanks must be used for each paragraph as LaTeX2e's \thanks
% was not built to handle multiple paragraphs
%

\author{
\IEEEauthorblockN{Lívia Gomes Costa Fonseca} \\
\and
\IEEEauthorblockN{Paulo Alvim Alvarenga}
}% <-this % stops a space

% \thanks{M. Shell was with the Department
% of Electrical and Computer Engineering, Georgia Institute of Technology, Atlanta,
% GA, 30332 USA e-mail: (see http://www.michaelshell.org/contact.html).}% <-this % stops a space
% \thanks{J. Doe and J. Doe are with Anonymous University.}% <-this % stops a space
% \thanks{Manuscript received April 19, 2005; revised August 26, 2015.}}

% note the % following the last \IEEEmembership and also \thanks - 
% these prevent an unwanted space from occurring between the last author name
% and the end of the author line. i.e., if you had this:
% 
% \author{....lastname \thanks{...} \thanks{...} }
%                     ^------------^------------^----Do not want these spaces!
%
% a space would be appended to the last name and could cause every name on that
% line to be shifted left slightly. This is one of those "LaTeX things". For
% instance, "\textbf{A} \textbf{B}" will typeset as "A B" not "AB". To get
% "AB" then you have to do: "\textbf{A}\textbf{B}"
% \thanks is no different in this regard, so shield the last } of each \thanks
% that ends a line with a % and do not let a space in before the next \thanks.
% Spaces after \IEEEmembership other than the last one are OK (and needed) as
% you are supposed to have spaces between the names. For what it is worth,
% this is a minor point as most people would not even notice if the said evil
% space somehow managed to creep in.


% The only time the second header will appear is for the odd numbered pages
% after the title page when using the twoside option.
% 
% *** Note that you probably will NOT want to include the author's ***
% *** name in the headers of peer review papers.                   ***
% You can use \ifCLASSOPTIONpeerreview for conditional compilation here if
% you desire.




% If you want to put a publisher's ID mark on the page you can do it like
% this:
%\IEEEpubid{0000--0000/00\$00.00~\copyright~2015 IEEE}
% Remember, if you use this you must call \IEEEpubidadjcol in the second
% column for its text to clear the IEEEpubid mark.



% use for special paper notices
%\IEEEspecialpapernotice{(Invited Paper)}




% make the title area
\maketitle

% As a general rule, do not put math, special symbols or citations
% in the abstract or keywords.
\begin{abstract}
Este trabalho apresenta a implementação do pipeline de detecção de campo e bola proposto por Nuruddin \textit{et al.} \cite{nuruddin2018} para robôs humanoides na plataforma EROS, aplicado a um conjunto de imagens capturadas pelas câmeras superior e inferior do robô NAO em competições RoboCup. O pipeline é composto por cinco etapas sequenciais: segmentação por cor no espaço HSV, morfologia matemática, detecção do contorno do campo, detecção de candidatos à bola por análise de blobs, e verificação dos candidatos por meio de descritores LBP (\textit{Local Binary Pattern}). Os resultados obtidos são comparados qualitativamente com os reportados no artigo original, e as limitações da abordagem clássica de segmentação por cor em cenários com iluminação variável e objetos da mesma cor do campo são discutidas.
\end{abstract}

% Note that keywords are not normally used for peerreview papers.
\begin{IEEEkeywords}
Detecção de bola, detecção de campo, segmentação HSV, morfologia matemática, LBP, RoboCup.
\end{IEEEkeywords}






% For peer review papers, you can put extra information on the cover
% page as needed:
% \ifCLASSOPTIONpeerreview
% \begin{center} \bfseries EDICS Category: 3-BBND \end{center}
% \fi
%
% For peerreview papers, this IEEEtran command inserts a page break and
% creates the second title. It will be ignored for other modes.
\IEEEpeerreviewmaketitle



% \section{Introduction}
% The very first letter is a 2 line initial drop letter followed
% by the rest of the first word in caps.
% 
% form to use if the first word consists of a single letter:
% \IEEEPARstart{A}{demo} file is ....
% 
% form to use if you need the single drop letter followed by
% normal text (unknown if ever used by the IEEE):
% \IEEEPARstart{A}{}demo file is ....
% 
% Some journals put the first two words in caps:
% \IEEEPARstart{T}{his demo} file is ....
% 
% Here we have the typical use of a "T" for an initial drop letter
% and "HIS" in caps to complete the first word.
% \IEEEPARstart{T}{his} demo file is intended to serve as a ``starter file''
% for IEEE journal papers produced under \LaTeX\ using
% IEEEtran.cls version 1.8b and later.
% You must have at least 2 lines in the paragraph with the drop letter
% (should never be an issue)
% I wish you the best of success.

% \hfill mds
 
% \hfill August 26, 2015

% \subsection{Subsection Heading Here}
% Subsection text here.

% needed in second column of first page if using \IEEEpubid
%\IEEEpubidadjcol

% \subsubsection{Subsubsection Heading Here}
% Subsubsection text here.


% An example of a floating figure using the graphicx package.
% Note that \label must occur AFTER (or within) \caption.
% For figures, \caption should occur after the \includegraphics.
% Note that IEEEtran v1.7 and later has special internal code that
% is designed to preserve the operation of \label within \caption
% even when the captionsoff option is in effect. However, because
% of issues like this, it may be the safest practice to put all your
% \label just after \caption rather than within \caption{}.
%
% Reminder: the "draftcls" or "draftclsnofoot", not "draft", class
% option should be used if it is desired that the figures are to be
% displayed while in draft mode.
%
%\begin{figure}[!t]
%\centering
%\includegraphics[width=2.5in]{myfigure}
% where an .eps filename suffix will be assumed under latex, 
% and a .pdf suffix will be assumed for pdflatex; or what has been declared
% via \DeclareGraphicsExtensions.
%\caption{Simulation results for the network.}
%\label{fig_sim}
%\end{figure}

% Note that the IEEE typically puts floats only at the top, even when this
% results in a large percentage of a column being occupied by floats.


% An example of a double column floating figure using two subfigures.
% (The subfig.sty package must be loaded for this to work.)
% The subfigure \label commands are set within each subfloat command,
% and the \label for the overall figure must come after \caption.
% \hfil is used as a separator to get equal spacing.
% Watch out that the combined width of all the subfigures on a 
% line do not exceed the text width or a line break will occur.
%
%\begin{figure*}[!t]
%\centering
%\subfloat[Case I]{\includegraphics[width=2.5in]{box}%
%\label{fig_first_case}}
%\hfil
%\subfloat[Case II]{\includegraphics[width=2.5in]{box}%
%\label{fig_second_case}}
%\caption{Simulation results for the network.}
%\label{fig_sim}
%\end{figure*}
%
% Note that often IEEE papers with subfigures do not employ subfigure
% captions (using the optional argument to \subfloat[]), but instead will
% reference/describe all of them (a), (b), etc., within the main caption.
% Be aware that for subfig.sty to generate the (a), (b), etc., subfigure
% labels, the optional argument to \subfloat must be present. If a
% subcaption is not desired, just leave its contents blank,
% e.g., \subfloat[].


% An example of a floating table. Note that, for IEEE style tables, the
% \caption command should come BEFORE the table and, given that table
% captions serve much like titles, are usually capitalized except for words
% such as a, an, and, as, at, but, by, for, in, nor, of, on, or, the, to
% and up, which are usually not capitalized unless they are the first or
% last word of the caption. Table text will default to \footnotesize as
% the IEEE normally uses this smaller font for tables.
% The \label must come after \caption as always.
%
%\begin{table}[!t]
%% increase table row spacing, adjust to taste
%\renewcommand{\arraystretch}{1.3}
% if using array.sty, it might be a good idea to tweak the value of
% \extrarowheight as needed to properly center the text within the cells
%\caption{An Example of a Table}
%\label{table_example}
%\centering
%% Some packages, such as MDW tools, offer better commands for making tables
%% than the plain LaTeX2e tabular which is used here.
%\begin{tabular}{|c||c|}
%\hline
%One & Two\\
%\hline
%Three & Four\\
%\hline
%\end{tabular}
%\end{table}


% Note that the IEEE does not put floats in the very first column
% - or typically anywhere on the first page for that matter. Also,
% in-text middle ("here") positioning is typically not used, but it
% is allowed and encouraged for Computer Society conferences (but
% not Computer Society journals). Most IEEE journals/conferences use
% top floats exclusively. 
% Note that, LaTeX2e, unlike IEEE journals/conferences, places
% footnotes above bottom floats. This can be corrected via the
% \fnbelowfloat command of the stfloats package.

\section{Introdução}

\IEEEPARstart{A}{} RoboCup é uma competição internacional de robótica e inteligência artificial criada em 1997 com o objetivo de estimular o desenvolvimento científico e tecnológico por meio de desafios práticos. Ela reúne estudantes, pesquisadores e profissionais de diversas áreas para projetar, construir e programar robôs capazes de executar tarefas de forma autônoma. A competição é organizada em diferentes categorias, como RoboCup Soccer (Futebol de Robôs), Rescue, @Home e Industrial, permitindo que equipes desenvolvam soluções para problemas reais em ambientes controlados e desafiadores.

Um dos objetivos de longo prazo da RoboCup é desenvolver, até meados do século XXI, uma equipe de robôs humanoides totalmente autônomos capaz de derrotar a equipe campeã mundial de futebol humano seguindo as regras oficiais da FIFA. Para isso, os robôs precisam perceber o ambiente por meio de câmeras, identificar a posição da bola, dos companheiros e dos adversários e executar ações como deslocamento, passe, marcação e chute ao gol.

Um desafio central nesse contexto é a percepção visual robusta. A partir de 2015, a RoboCup passou a utilizar bola e gol inteiramente brancos, eliminando marcações coloridas que facilitavam a detecção por cor. Esse trabalho implementa e avalia o pipeline proposto por Nuruddin \textit{et al.} \cite{nuruddin2018}, que combina segmentação por cor, morfologia matemática e classificador em cascata com descritores LBP para detectar o campo e a bola em condições adversas.

\section{Metodologia}

O pipeline implementado segue as cinco etapas descritas em \cite{nuruddin2018}. As imagens utilizadas para avaliação são apresentadas na Fig.~\ref{fig:entrada_grid}.

\subsection{Dataset}

As imagens utilizadas neste trabalho prov\^{e}m do \textit{RoboCup SPL Instance Segmentation Dataset} \cite{naodevils2019}, disponibilizado pela equipe NaoDevils no Kaggle. O dataset cont\'em 1.179 imagens segmentadas manualmente e 8.821 anotadas automaticamente, capturadas pela c\^{a}mera superior do rob\^{o} NAO durante a RoboCup~2019 e a GermanOpen~2019, ambas na Standard Platform League (SPL). As categorias anotadas incluem: bola, rob\^{o}, linha, c\'irculo central, gol e marcador de p\^{e}nalti. As imagens foram utilizadas tanto para testar o pipeline (etapas~1 a~4) quanto para treinar o classificador LBP da etapa~5, aproveitando as anota\c{c}\~{o}es de bola no formato COCO.

A Tabela~\ref{tab:imagens} lista as 12 imagens disponíveis e seu conteúdo. Elas foram divididas em dois grupos: imagens sem bola, usadas para avaliar a detecção do campo, e imagens com bola, usadas para avaliar o pipeline completo.

\begin{table}[!t]
\renewcommand{\arraystretch}{1.3}
\caption{Imagens de Entrada Utilizadas}
\label{tab:imagens}
\centering
\begin{tabular}{|l|c|l|}
\hline
\textbf{Arquivo} & \textbf{Bola?} & \textbf{Característica principal} \\
\hline
\texttt{campo-normal.png}   & Não & Iluminação uniforme, campo limpo \\
\texttt{campo-escuro.png}   & Não & Iluminação reduzida \\
\texttt{campo-gol.png}      & Não & Gol visível ao fundo \\
\texttt{campo-vazio.png}    & Não & Campo vazio, sem bola \\
\texttt{campo-1.png}        & Não & Vista alternativa do campo \\
\hline
\texttt{bola-boa.png}       & Sim & Bola nítida, posição central \\
\texttt{bola-boa-2.png}     & Sim & Bola nítida, ângulo diferente \\
\texttt{bola-na-linha.png}  & Sim & Bola sobre linha branca \\
\texttt{bola-com-robo.png}  & Sim & Bola próxima a robô \\
\texttt{bola-borrado.png}   & Sim & Bola com motion blur \\
\texttt{penalty-boad.png}   & Sim & Bola em situação de pênalti \\
\hline
\end{tabular}
\end{table}

\begin{figure*}[!t]
\centering
\includegraphics[width=6.8in]{figuras/dataset_imagens_entrada.png}
\caption{As 12 imagens de entrada utilizadas neste trabalho, capturadas pela câmera superior do robô NAO em competições RoboCup~2019. Bordas \textbf{\textcolor[rgb]{0.12,0.47,0.71}{azuis}}: imagens sem bola, usadas para avaliar a detecção do campo; bordas \textbf{\textcolor[rgb]{0.17,0.63,0.17}{verdes}}: imagens com bola, usadas para avaliar o pipeline completo.}
\label{fig:entrada_grid}
\end{figure*}

\subsection{Etapa 1 — Segmentação por Cor no Espaço HSV}

O primeiro passo é isolar a região do campo na imagem. O campo é composto por grama sintética verde, tornando a segmentação por cor uma abordagem direta. O espaço de cor HSV (\textit{Hue, Saturation, Value}) é preferível ao RGB porque separa a informação de matiz (cor) da luminosidade, tornando a segmentação mais robusta a variações de iluminação.

Foi aplicado o operador \texttt{cv2.inRange} com os limites $H \in [30, 90]$, $S \in [40, 255]$ e $V \in [40, 255]$, conforme proposto no paper. O resultado é uma máscara binária onde os pixels pertencentes ao campo recebem valor 255 e os demais, zero.

\subsection{Etapa 2 — Morfologia Matemática}

A máscara obtida na etapa anterior apresenta ruídos: pequenos buracos no interior do campo causados pelas linhas brancas, e pequenos pontos verdes fora do campo provenientes de placas publicitárias e da torcida. Para limpar a máscara, foram aplicadas as seguintes operações morfológicas com elemento estruturante elíptico de tamanho $5 \times 5$:

\begin{enumerate}
    \item \textbf{Abertura} (erosão seguida de dilatação): remove componentes conexas menores que o elemento estruturante, eliminando ruídos externos ao campo.
    \item \textbf{Fechamento} (dilatação seguida de erosão): preenche buracos menores que o elemento estruturante, conectando regiões do campo separadas pelas linhas brancas.
    \item \textbf{Abertura final}: remove ilhas verdes remanescentes fora do campo principal.
\end{enumerate}

O elemento estruturante elíptico foi escolhido por aproximar melhor a forma circular da bola e as bordas curvas do campo em comparação a um elemento quadrado.

A Fig.~\ref{fig:morfologia_normal} ilustra o efeito de cada operação na imagem \texttt{campo-normal.png}: a abertura elimina pontos verdes isolados fora do campo, o fechamento preenche os buracos causados pelas linhas brancas e a abertura final produz uma máscara limpa e coesa. A Fig.~\ref{fig:morfologia_escuro} mostra o mesmo processo em \texttt{campo-escuro.png}, onde a iluminação reduzida fragmenta a máscara original em várias regiões, evidenciando a maior dificuldade do método em condições adversas.

\begin{figure}[!t]
\centering
\includegraphics[width=3.3in]{figuras/resultado_etapa2_campo-normal_png.png}
\caption{Morfologia aplicada a \texttt{campo-normal.png}: (a)~máscara original (HSV), (b)~após abertura, (c)~após fechamento, (d)~máscara final. A progressão mostra a remoção de ruídos externos e o preenchimento das linhas brancas.}
\label{fig:morfologia_normal}
\end{figure}

\begin{figure}[!t]
\centering
\includegraphics[width=3.3in]{figuras/resultado_etapa2_campo-escuro_png.png}
\caption{Morfologia aplicada a \texttt{campo-escuro.png}: a iluminação baixa produz uma máscara inicial fragmentada (a), e as operações morfológicas conseguem unificar parte das regiões (b--d), mas a máscara final ainda apresenta lacunas, afetando a estimativa do contorno do campo na etapa seguinte.}
\label{fig:morfologia_escuro}
\end{figure}

\subsection{Etapa 3 — Detecção do Campo}

Com a máscara limpa, foi aplicado \texttt{cv2.findContours} para identificar os contornos presentes. O maior contorno por área é selecionado como o contorno do campo. Em seguida, o \textit{convex hull} do contorno é calculado para obter uma estimativa suavizada da borda do campo.

A \textit{linha do horizonte} é determinada coluna a coluna: para cada coluna $x$ da imagem, registra-se o menor valor de $y$ onde a máscara é positiva. Essa linha delimita a fronteira superior do campo e é usada para restringir a busca por candidatos à bola apenas à região inferior (dentro do campo).

A \textit{Region of Interest} (ROI) é definida como o \textit{bounding box} do maior contorno, reduzindo a área de busca das etapas seguintes.

\subsection{Etapa 4 — Detecção de Candidatos à Bola}

Dentro da ROI do campo, busca-se por regiões que não são verdes e que possuem características geométricas compatíveis com uma bola. A bola do RoboCup é branca com padrão preto, portanto:

\begin{enumerate}
    \item A máscara do campo é invertida, destacando o que não é verde dentro da ROI.
    \item Um limiar de Otsu é aplicado à imagem em escala de cinza para isolar regiões claras (bola branca).
    \item As duas máscaras são combinadas por operação \texttt{AND}: regiões claras e não-verdes dentro da ROI.
    \item Os contornos resultantes são filtrados por área ($80 \leq A \leq 8000$ pixels$^2$) e \textbf{circularidade} $C = \frac{4\pi A}{P^2}$, com $C \geq 0{,}35$, onde $P$ é o perímetro do contorno.
\end{enumerate}

Os candidatos são ordenados por circularidade decrescente, de forma que o mais provável de ser a bola apareça primeiro.

\subsection{Etapa 5 — Classificador com Descritores LBP}

O paper utiliza um classificador em cascata treinado com descritores LBP (\textit{Local Binary Pattern}) para validar os candidatos. O LBP é calculado comparando cada pixel com seus $P$ vizinhos em um raio $R$: se o vizinho for maior ou igual ao pixel central, atribui-se 1; caso contrário, 0. O resultado é um número binário de $P$ bits cujo histograma descreve a textura local da região.

O LBP foi escolhido pelo paper por ser computacionalmente leve, adequado ao hardware embarcado dos robôs. Em nossa implementação, foi treinado um classificador SVM com kernel RBF utilizando histogramas LBP uniformes ($P=8$, $R=1$, 10 bins) extraídos de patches do dataset NaoDevils \cite{naodevils2019}, composto por 296 anotações de bola (amostras positivas) e o triplo em amostras negativas extraídas aleatoriamente das mesmas imagens. O classificador é aplicado a cada patch candidato redimensionado para $32 \times 32$ pixels; um candidato é aceito como bola se a probabilidade estimada pelo SVM supera o limiar de $0{,}45$.

A Fig.~\ref{fig:etapa5_resultado} ilustra o processo de classificação para a imagem \texttt{bola-com-robo.png}: os candidatos são numerados e exibidos com seus respectivos scores, e o mapa LBP da imagem inteira é sobreposto com os círculos de detecção. Para cada candidato, o patch recortado e seu histograma LBP são exibidos individualmente (Fig.~\ref{fig:etapa5_patches}).

\begin{figure}[!t]
\centering
\includegraphics[width=3.3in]{figuras/resultado_etapa5_bola-com-robo_png.png}
\caption{Visão geral da Etapa~5 para \texttt{bola-com-robo.png}: (1)~imagem original, (2)~candidatos numerados, (3)~mapa LBP com candidatos sobrepostos, (4)~resultado final (verde=bola, vermelho=rejeitado), (5)~histograma LBP global, (6)~score de cada candidato com linha do limiar.}
\label{fig:etapa5_resultado}
\end{figure}

\begin{figure}[!t]
\centering
\includegraphics[width=3.3in]{figuras/resultado_etapa5_patches_bola-com-robo_png.png}
\caption{Detalhe LBP por candidato em \texttt{bola-com-robo.png}: para cada candidato são exibidos o patch recortado (escala de cinza), o mapa LBP do patch e o histograma LBP correspondente. Candidatos em verde foram classificados como bola; em vermelho, rejeitados.}
\label{fig:etapa5_patches}
\end{figure}

\section{Resultados}

\subsection{Detecção do Campo}

Os resultados da detecção do campo foram avaliados nas cinco imagens sem bola disponíveis: \texttt{campo-normal.png}, \texttt{campo-escuro.png}, \texttt{campo-gol.png}, \texttt{campo-vazio.png} e \texttt{campo-1.png}.

\begin{table}[!t]
\renewcommand{\arraystretch}{1.3}
\caption{Resultados da Detecção do Campo}
\label{tab:campo}
\centering
\begin{tabular}{|l|c|c|c|}
\hline
\textbf{Imagem} & \textbf{Cobertura (\%)} & \textbf{Contornos} & \textbf{Campo detectado?} \\
\hline
campo-normal     & 88,3 & 1 & Sim \\
campo-escuro     & 63,3 & 5 & Sim (parcial) \\
campo-gol        & 69,1 & 7 & Sim (parcial) \\
campo-sem-bola   & 69,9 & 9 & Sim (parcial) \\
\hline
\end{tabular}
\end{table}

Na imagem \texttt{campo-normal.png}, sob iluminação uniforme e sem elementos externos no enquadramento, o algoritmo detectou corretamente o campo com 88,3\% de cobertura e apenas um contorno, indicando boa segmentação. Nas demais imagens, a presença de placas publicitárias verdes, torcida e variação de iluminação aumentou o número de contornos e reduziu a cobertura estimada. Em particular, a imagem \texttt{campo-escuro.png} apresentou fragmentação significativa da máscara, com 5 contornos distintos, evidenciando a sensibilidade do método a variações de iluminação.

\begin{figure}[!t]
\centering
\subfloat[Segmentação HSV]{\includegraphics[width=1.55in]{figuras/resultado_etapa1_campo-normal_png.png}}
\hfill
\subfloat[Morfologia + Contorno]{\includegraphics[width=1.55in]{figuras/resultado_etapa3_campo-normal_png.png}}
\caption{Etapas 1 e 3 aplicadas à imagem \texttt{campo-normal.png}.}
\label{fig:campo_normal}
\end{figure}

Um problema recorrente observado foi que a \textit{linha do horizonte} sobe para além da borda real do campo em imagens onde objetos verdes (placas, roupas da torcida) aparecem na região superior. Como a linha do horizonte é calculada pelo primeiro pixel verde de cada coluna, qualquer pixel verde espúrio no topo da imagem puxa a estimativa para cima, fazendo com que a ROI inclua região fora do campo e exclua parte da borda inferior. Esse comportamento é ilustrado na Fig.~\ref{fig:campo_escuro}, onde a iluminação reduzida faz com que pixels verdes do campo se percam na segmentação e fragmentos verdes acima do campo (torcida, placas) distorcem a estimativa do horizonte para cima.

\begin{figure}[!t]
\centering
\includegraphics[width=3.3in]{figuras/resultado_etapa3_campo-escuro_png.png}
\caption{Resultado da detecção de campo em \texttt{campo-escuro.png}. A linha do horizonte estimada (vermelho) situa-se acima da borda real do campo devido a pixels verdes espúrios na região superior da imagem. A ROI resultante inclui área fora do campo e pode excluir a borda inferior, onde a bola costuma estar posicionada.}
\label{fig:campo_escuro}
\end{figure}

É importante notar que um campo \textbf{maior} (maior cobertura da imagem) é geralmente mais favorável para a detecção da bola do que um campo menor. Quando a ROI é ampla, a busca por candidatos abrange mais área útil e tolera melhor erros na estimativa da linha do horizonte. Por outro lado, quando a cobertura do campo é pequena — como ocorre em \texttt{campo-escuro.png} com 63,3\% — a ROI fica reduzida e uma bola posicionada na borda do campo pode ficar parcialmente fora da região de busca, resultando em falha de detecção. Esse trade-off é uma limitação inerente à abordagem de restrição por ROI adotada pelo pipeline.

\subsection{Detecção da Bola}

A detecção da bola foi avaliada nas sete imagens que contêm bola: \texttt{bola-boa.png}, \texttt{bola-boa-2.png}, \texttt{bola-borrado.png}, \texttt{bola-com-robo.png}, \texttt{bola-na-linha.png}, \texttt{PENALTY-BORRADO.png} e \texttt{penalty-boad.png}.

\begin{table}[!t]
\renewcommand{\arraystretch}{1.3}
\caption{Resultados da Detecção da Bola}
\label{tab:bola}
\centering
\begin{tabular}{|l|c|c|}
\hline
\textbf{Imagem} & \textbf{Candidatos} & \textbf{Bola detectada?} \\
\hline
bola-boa          & $\geq 1$ & Sim \\
bola-boa-2        & $\geq 1$ & Sim \\
bola-na-linha     & $\geq 1$ & Parcial \\
bola-com-robo     & $\geq 1$ & Parcial \\
bola-borrado      & $\geq 1$ & Não \\
PENALTY-BORRADO   & $\geq 1$ & Não \\
penalty-boad      & $\geq 1$ & Sim \\
\hline
\end{tabular}
\end{table}

Nas imagens com bola nítida e bem posicionada dentro do campo (\texttt{bola-boa.png} e \texttt{bola-boa-2.png}), o pipeline identificou candidatos com circularidade elevada e o classificador LBP atribuiu score suficiente para confirmar a detecção. Nas imagens com bola parcialmente borrada (\texttt{bola-borrado.png}, \texttt{PENALTY-BORRADO.png}), o borramento reduz o contraste e deforma o contorno, diminuindo a circularidade abaixo do limiar e impedindo a detecção.

Um falso positivo relevante foi observado de forma consistente na imagem \texttt{bola-com-robo.png}: o pé do robô, que possui coloração preto-e-branca similar à bola, gerou candidatos com score elevado (acima de $0{,}9$). Esse comportamento revela uma limitação do treinamento: como o dataset de amostras negativas é composto por patches aleatórios do campo (majoritariamente grama), o classificador aprende a distinguir bola de grama, mas não de outros objetos bicolores como os pés e canelas dos robôs. Um treinamento mais robusto incluiria amostras negativas explícitas de membros de robôs (\textit{hard negatives}), reduzindo essa categoria de falso positivo. A Fig.~\ref{fig:etapa5_resumo} apresenta o resumo comparativo de todas as imagens processadas, mostrando a quantidade de candidatos e bolas detectadas por imagem, além da distribuição dos scores.

\begin{figure}[!t]
\centering
\includegraphics[width=3.3in]{figuras/resultado_etapa5_resumo_comparativo.png}
\caption{Resumo comparativo da Etapa~5: quantidade de candidatos e bolas detectadas por imagem (esquerda) e score máximo e médio por imagem com linha do limiar (direita).}
\label{fig:etapa5_resumo}
\end{figure}

\subsection{Comparação com o Paper Original}

O paper \cite{nuruddin2018} reporta resultados experimentais em imagens capturadas pela câmera frontal do robô EROS em competições indonesianas de robótica humanoide. O artigo não disponibiliza métricas numéricas precisas (como \textit{precision} e \textit{recall}), reportando os resultados qualitativamente como ``encorajadores'' e descrevendo que a combinação de segmentação por cor, morfologia e classificador em cascata reduziu significativamente os falsos positivos em comparação ao uso isolado da segmentação por cor.

A Tabela~\ref{tab:comparacao} resume as principais diferenças entre o cenário do paper e o nosso.

\begin{table*}[!t]
\renewcommand{\arraystretch}{1.3}
\caption{Comparação com o Paper Original}
\label{tab:comparacao}
\centering
\begin{tabular}{|l|p{2.6in}|p{2.6in}|}
\hline
\textbf{Aspecto} & \textbf{Paper original} & \textbf{Nossa implementação} \\
\hline
Câmera utilizada & Frontal embarcada (robô EROS) & Superior e inferior do robô NAO \\
\hline
Classificador & Cascade Classifier treinado com LBP & SVM com kernel RBF treinado com histogramas LBP (P=8, R=1) \\
\hline
Espaço de cor & YUV & HSV \\
\hline
Campo detectado & Sim (condições controladas) & Sim (parcial em imagens com iluminação variável ou elementos externos) \\
\hline
Bola detectada & Sim (resultado encorajador) & Sim em casos simples; falsa detecção de pés de robôs; falha em borramentos \\
\hline
\end{tabular}
\end{table*}

As principais diferenças de desempenho decorrem de quatro fatores. Primeiro, a câmera: o paper utiliza câmera frontal embarcada no robô EROS, que enquadra o campo de forma mais direta, com menos elementos externos (torcida, placas) no quadro. Nossas imagens são capturadas pela câmera superior do robô NAO (SPL) — perspectiva de cima que, dependendo da orientação do robô, inclui frequentemente a torcida e as bordas do campo, comprometendo a segmentação por cor e a estimativa do horizonte.

Segundo, o domínio do dataset: o classificador LBP foi treinado com imagens do dataset NaoDevils \cite{naodevils2019}, coletadas pela câmera superior do NAO na RoboCup~2019 e GermanOpen~2019. O paper EROS treina e testa no mesmo domínio (câmera frontal, robô EROS, competições indonésias). Essa diferença de domínio — câmera, ângulo, iluminação e tipo de competição distintos — limita a transferência direta dos resultados e é uma fonte adicional de falsos positivos.

Terceiro, o classificador: o paper treina um Cascade Classifier offline com amostras do dataset específico da equipe EROS, obtendo um modelo ajustado ao seu domínio. Nossa implementação treina um SVM com histogramas LBP em um dataset público (NaoDevils), o que é funcionalmente equivalente mas com menor diversidade de cenários negativos — daí a ocorrência de falsos positivos em membros de robôs.

Quarto, o espaço de cor: o paper usa YUV, enquanto usamos HSV. Ambos separam luminância de crominância, mas com diferentes sensibilidades. O YUV é nativamente produzido por algumas câmeras, evitando conversão; o HSV é mais intuitivo para definição de faixas de cor por humanos. Na prática, a escolha do espaço de cor é menos determinante do que a qualidade da calibração dos limiares.

\section{Conclusão}

Este trabalho implementou e avaliou o pipeline de detecção de campo e bola proposto por Nuruddin \textit{et al.} \cite{nuruddin2018} em um conjunto de imagens capturadas pelas câmeras superior e inferior do robô NAO em competições RoboCup. Os resultados mostraram que a abordagem funciona bem em condições controladas — iluminação uniforme, câmera frontal, sem elementos externos ao campo no quadro — mas apresenta limitações importantes em cenários mais realistas.

A principal limitação observada é a sensibilidade da segmentação por cor a variações de iluminação e à presença de objetos da mesma cor do campo fora dele. Isso afeta diretamente a qualidade da máscara morfológica e, consequentemente, a estimativa da linha do horizonte e da ROI. Imagens capturadas pelas câmeras superior e inferior do NAO agravam esse problema por incluir frequentemente a torcida, placas publicitárias e outros elementos externos ao campo.

A detecção da bola funciona de forma razoável para bolas nítidas e bem posicionadas, mas falha sistematicamente em imagens borradas, onde a deformação do contorno reduz a circularidade abaixo do limiar de detecção. O classificador SVM com LBP, embora treinado com dados reais de competições RoboCup, apresenta falsos positivos em objetos bicolores como os pés dos robôs, pois as amostras negativas de treinamento são dominadas por grama e não cobrem adequadamente essa categoria.

Em suma, a abordagem proposta pelo paper é viável para hardware embarcado com recursos limitados, mas exige calibração cuidadosa dos parâmetros HSV para cada ambiente de iluminação e um classificador treinado com amostras do domínio de aplicação específico.


% Templates úteis 

% Imagem

% \begin{figure}[H]
%     \centering
%     \includegraphics[width=2in]{imagens/}
%     \caption{Legenda}
%     \label{}
% \end{figure}

%  Imagem com duas subfiguras

% \begin{figure}[H]
%     \centering
%     \subfloat[Legenda 1]{\includegraphics[width=1.5in]{imagens/}%
%     }
%     \hfill
%     \subfloat[Legenda 2]{\includegraphics[width=1.5in]{imagens/}%
%     }
%     \caption{Legenda}
%     \label{}
% \end{figure}

% Algoritmo

% \begin{algorithm} [H]
%     \caption{Legenda}
%     \begin{algorithmic}
%     \STATE instrução
%     \FOR{for}
%         \STATE instrução
%     \ENDFOR
%     \end{algorithmic}
% \end{algorithm}

\appendix
O código do trabalho pode ser acessado pelo link do github: \url{https://github.com/alvimpaulo/IPI-2026.1}





% if have a single appendix:
%\appendix[Proof of the Zonklar Equations]
% or
%\appendix  % for no appendix heading
% do not use \section anymore after \appendix, only \section*
% is possibly needed

% use appendices with more than one appendix
% then use \section to start each appendix
% you must declare a \section before using any
% \subsection or using \label (\appendices by itself
% starts a section numbered zero.)
%


% \appendices
% \section{Proof of the First Zonklar Equation}
% Appendix one text goes here.

% you can choose not to have a title for an appendix
% if you want by leaving the argument blank
% \section{}
% Appendix two text goes here.


% use section* for acknowledgment
% \section*{Acknowledgment}


% The authors would like to thank...


% Can use something like this to put references on a page
% by themselves when using endfloat and the captionsoff option.
% \ifCLASSOPTIONcaptionsoff
%   \newpage
% \fi



% trigger a \newpage just before the given reference
% number - used to balance the columns on the last page
% adjust value as needed - may need to be readjusted if
% the document is modified later
%\IEEEtriggeratref{8}
% The "triggered" command can be changed if desired:
%\IEEEtriggercmd{\enlargethispage{-5in}}

% references section

% can use a bibliography generated by BibTeX as a .bbl file
% BibTeX documentation can be easily obtained at:
% http://mirror.ctan.org/biblio/bibtex/contrib/doc/
% The IEEEtran BibTeX style support page is at:
% http://www.michaelshell.org/tex/ieeetran/bibtex/
\begin{thebibliography}{2}

\bibitem{nuruddin2018}
M.~N. Sudin, S.~R. S.~Abdullah, and M.~F. Nasudin,
``Improving Field and Ball Detector for Humanoid Robot Soccer EROS Platform,\''
in \emph{Proc. IEEE Int. Conf. Electrical Systems and Power Engineering (ELECSYM)},
2018, pp.~1--6. doi: 10.1109/elecsym.2018.8615506

\bibitem{naodevils2019}
P.~Broemmel,
``RoboCup SPL Instance Segmentation Dataset,''
\textit{Kaggle}, 2020.
[Online]. Available: \url{https://www.kaggle.com/datasets/pietbroemmel/naodevils-segmentation-upper-camera}

\end{thebibliography}

% biography section
% 
% If you have an EPS/PDF photo (graphicx package needed) extra braces are
% needed around the contents of the optional argument to biography to prevent
% the LaTeX parser from getting confused when it sees the complicated
% \includegraphics command within an optional argument. (You could create
% your own custom macro containing the \includegraphics command to make things
% simpler here.)
%\begin{IEEEbiography}[{\includegraphics[width=1in,height=1.25in,clip,keepaspectratio]{mshell}}]{Michael Shell}
% or if you just want to reserve a space for a photo:

% \begin{IEEEbiography}{Michael Shell}
% Biography text here.
% \end{IEEEbiography}

% if you will not have a photo at all:
% \begin{IEEEbiographynophoto}{John Doe}
% Biography text here.
% \end{IEEEbiographynophoto}

% insert where needed to balance the two columns on the last page with
% biographies
%\newpage

% \begin{IEEEbiographynophoto}{Jane Doe}
% Biography text here.
% \end{IEEEbiographynophoto}

% You can push biographies down or up by placing
% a \vfill before or after them. The appropriate
% use of \vfill depends on what kind of text is
% on the last page and whether or not the columns
% are being equalized.

%\vfill

% Can be used to pull up biographies so that the bottom of the last one
% is flush with the other column.
%\enlargethispage{-5in}



% that's all folks
\end{document}


